\section{Conclusion}

We introduce \oursys to characterize the capability boundary of LLM-based Triton kernel generation through category-aware evaluation across 176 tasks on six GPUs. Three insights emerge.
The correctness boundary is category-structured and driven by non-local semantic coordination.
Iterative refinement is repair-biased: it expands the feasible set but introduces weaker candidates, the edit distribution is dominated by local fixes.
Performance validity remains a distinct and unsolved challenge.

Together, these results indicate that the capability boundary of current LLM-based kernel generation is not a single wall but a sequence of distinct barriers - compilability, semantic correctness, hardware efficiency and performance portability - each requiring different mechanisms to clear.
Prompt engineering and iterative refinement are well-suited to compilability but structurally insufficient for the rest.
Progress will likely require mechanisms for reasoning about global tensor contracts and parallel reduction semantics, training signals that reward numerical fidelity rather than surface resemblance, and efficiency-aware generation with explicit hardware cost feedback.
We release \oursys together with the iteration-level transition pairs collected during evaluation to support research in these directions.
